\documentclass{beamer}
\usetheme{Madrid}

\usepackage[utf8]{inputenc} % solo si usas pdflatex
\usepackage[T1]{fontenc}    % solo si usas pdflatex

\usepackage{amsmath,amssymb}
\usepackage{microtype}
\usepackage{cancel}         % si lo usas
\usepackage{fontawesome}    % si lo usas
% \usepackage{xypic}        % solo si lo usas
% \usepackage{multicol}     % normalmente no; usar columns de beamer
% \usepackage{amsthm}       % solo si defines teoremas propios
% \usepackage{scrextend}    % solo si lo usas

\usepackage{graphicx}
\graphicspath{{images/}}

\hypersetup{pdfborder={0 0 0}}


\title{Procesamiento distribuido con Spark}
\author{Rodrigo Blanco Betes. Alfa Formación}
\date{Enero del 2026}

\begin{document}

\begin{frame}
\titlepage
\end{frame}

\begin{frame}
\tableofcontents
\end{frame}

% ------------------------------------------------------------
% Sección 1
% ------------------------------------------------------------
\section{Uso de Spark como framework de procesamiento de la información en sistemas distribuidos}

% -------------------------

\begin{frame}{Principales retos en los sistemas distribuidos}
\begin{itemize}
  \item \textbf{Latencia y red}: comunicación entre nodos más lenta que la memoria local; penaliza algoritmos con muchas rondas.
  \item \textbf{Fallos parciales}: nodos que caen, discos que fallan, timeouts; el sistema debe \emph{recuperar} sin parar todo.
  \item \textbf{Consistencia y coordinación}: sincronización, locks, consenso; evitar cuellos de botella por coordinación excesiva.
  \item \textbf{Distribución de datos}: particionado, \emph{skew} (particiones desbalanceadas), localización (data locality).
  \item \textbf{Reproducibilidad y determinismo}: reintentos pueden duplicar efectos si no hay cuidado (\emph{exactly-once} vs \emph{at-least-once}).
  \item \textbf{Gestión de recursos}: memoria, CPU, disco, I/O; contención cuando varios jobs compiten.
\end{itemize}
\end{frame}

\begin{frame}{Retos típicos que aparecen en la práctica}
\begin{itemize}
  \item \textbf{Shuffles caros}: joins/aggregations provocan intercambio masivo de datos por red.
  \item \textbf{Stragglers}: unas pocas tareas lentas alargan toda la etapa (\emph{stage}).
  \item \textbf{Observabilidad}: entender “qué pasa” requiere métricas, UI, logs, y explicar planes de ejecución.
  \item \textbf{Coste}: más nodos no siempre implica más rendimiento (ley de rendimientos decrecientes).
\end{itemize}
\end{frame}

% -------------------------
\begin{frame}{¿Por qué Spark? Ventajas clave}
\begin{itemize}
  \item \textbf{Procesamiento en memoria}: acelera workloads iterativos (ML, graph, ETL encadenado).
  \item \textbf{Modelo unificado}: batch, streaming, SQL, ML (ecosistema integrado).
  \item \textbf{Tolerancia a fallos}: recomputación a partir de \emph{lineage} (DAG) de transformaciones.
  \item \textbf{Optimización automática}: Catalyst (planes lógicos/físicos) y Tungsten (ejecución eficiente).
  \item \textbf{APIs de alto nivel}: DataFrames/Datasets (tipado en Scala), SQL, funciones built-in.
  \item \textbf{Escalabilidad}: desde local a clúster (YARN, Kubernetes, Standalone) sin reescribir el código.
\end{itemize}
\end{frame}

\begin{frame}{Conceptos Spark que conectan con “beneficios”}
\begin{itemize}
  \item \textbf{DAG de ejecución}: el job se divide en \textbf{stages} (por barreras de shuffle) y \textbf{tasks}.
  \item \textbf{Transformations vs Actions}:
  \begin{itemize}
    \item \emph{Transformations} son \textbf{perezosas} (lazy): construyen el plan.
    \item \emph{Actions} disparan la ejecución: materializan el resultado.
  \end{itemize}
  \item \textbf{Persistencia/caching}: control explícito de reutilización para evitar recomputaciones.
  \item \textbf{Particionado}: clave para rendimiento (paralelismo) y para evitar \emph{skew}.
\end{itemize}
\end{frame}

\begin{frame}[fragile]{Ejemplo mínimo: lazy + action (Scala)}
\begin{verbatim}
val df = spark.read.parquet("/data/events")

val cleaned = df
  .filter($"country".isNotNull)
  .withColumn("day", to_date($"ts"))

val agg = cleaned.groupBy($"day").count()

agg.show()   // Action: aquí se ejecuta el plan
\end{verbatim}
\end{frame}

\begin{frame}{Cuándo Spark brilla (y cuándo no)}
\begin{itemize}
  \item \textbf{Brilla}: ETL grande, joins/aggregations a escala, pipelines, análisis exploratorio, streaming.
  \item \textbf{No siempre ideal}: latencia ultra-baja (ms), OLTP, tareas pequeñas donde el overhead domina.
  \item \textbf{Regla práctica}: Spark es una “fábrica” para datos; no un motor transaccional.
\end{itemize}
\end{frame}

% -------------------------


\begin{frame}{Fundamentos de programación funcional (FP)}
\begin{itemize}
  \item \textbf{Funciones como valores}: se pasan como parámetros y se retornan.
  \item \textbf{Inmutabilidad}: evita efectos colaterales y simplifica concurrencia.
  \item \textbf{Funciones puras}: mismo input \(\Rightarrow\) mismo output; facilita pruebas y reintentos.
  \item \textbf{Transformaciones sobre colecciones}: \texttt{map}, \texttt{filter}, \texttt{flatMap}, \texttt{reduce}.
  \item \textbf{Composición}: construir pipelines encadenando transformaciones.
\end{itemize}
\end{frame}

\begin{frame}{FP y sistemas distribuidos: por qué encaja tan bien}
\begin{itemize}
  \item \textbf{Menos estado compartido} \(\Rightarrow\) menos coordinación entre nodos.
  \item \textbf{Reintentos seguros} si las funciones son puras (muy alineado con tolerancia a fallos).
  \item \textbf{Pipelines declarativos}: describes \emph{qué} hacer; el motor decide \emph{cómo} ejecutarlo.
  \item \textbf{Paralelismo natural}: operaciones por elemento/partición se distribuyen fácilmente.
\end{itemize}
\end{frame}

\begin{frame}[fragile]{Ejemplo FP: de listas a RDD/DataFrame (Scala)}
\begin{verbatim}
// Colección local (FP clásica)
val xs = List(1,2,3,4,5)
val ys = xs.filter(_ % 2 == 0).map(_ * 10)

// Idea similar en Spark (distribuido)
val rdd = sc.parallelize(xs)
val out = rdd.filter(_ % 2 == 0).map(_ * 10)
\end{verbatim}
\end{frame}

\begin{frame}{Idea clave para Spark: "declarar transformaciones"}
\begin{itemize}
  \item En Spark, muchas operaciones son \textbf{transformaciones} (lazy) sobre datos distribuidos.
  \item La FP ayuda a pensar en términos de:
  \begin{itemize}
    \item \textbf{pipeline} de pasos
    \item \textbf{datos inmutables}
    \item \textbf{funciones} que transforman datos
  \end{itemize}
  \item Resultado: código más \textbf{razonable}, \textbf{testeable} y \textbf{escalable}.
\end{itemize}
\end{frame}

% ------------------------------------------------------------
% Sección 2
% ------------------------------------------------------------
\section{Inmersión en la estructura de datos principal de Spark: RDD -- Resilient Distributed Dataset}

\begin{frame}{¿Qué es un RDD?}
\begin{itemize}
  \item \textbf{RDD (Resilient Distributed Dataset)}: colección \textbf{distribuida} e \textbf{inmutable} de elementos.
  \item \textbf{Particionado}: los datos se dividen en particiones procesadas en paralelo.
  \item \textbf{Resilient}: tolerancia a fallos mediante \textbf{lineage} (DAG de transformaciones).
  \item Nivel de API más \textbf{bajo} que DataFrames/Datasets: más control, menos optimización automática.
\end{itemize}
\end{frame}

\begin{frame}{Tipos de operaciones sobre RDDs}
\begin{itemize}
  \item \textbf{Transformations} (lazy): construyen el plan, no ejecutan.
  \begin{itemize}
    \item \texttt{map}, \texttt{flatMap}, \texttt{filter}
    \item \texttt{distinct}, \texttt{union}, \texttt{intersection}
    \item \texttt{groupByKey}, \texttt{reduceByKey}, \texttt{join}
  \end{itemize}
  \item \textbf{Actions}: disparan la ejecución y devuelven resultado o escriben salida.
  \begin{itemize}
    \item \texttt{count}, \texttt{collect}, \texttt{take}, \texttt{first}
    \item \texttt{reduce}, \texttt{saveAsTextFile}
  \end{itemize}
  \item \textbf{Persistencia}: \texttt{cache}/\texttt{persist} para reutilizar resultados entre acciones.
\end{itemize}
\end{frame}

\begin{frame}{Operaciones ``narrow'' vs ``wide'' (impacto en rendimiento)}
\begin{itemize}
  \item \textbf{Narrow transformations}: cada partición de salida depende de una sola partición de entrada.
  \begin{itemize}
    \item Ejemplos: \texttt{map}, \texttt{filter}
    \item Ventaja: se ejecutan \textbf{sin shuffle}, suelen ser más rápidas.
  \end{itemize}
  \item \textbf{Wide transformations}: la salida depende de múltiples particiones de entrada.
  \begin{itemize}
    \item Ejemplos: \texttt{reduceByKey}, \texttt{groupByKey}, \texttt{join}
    \item Implican \textbf{shuffle}: intercambio de datos por red (coste alto).
  \end{itemize}
\end{itemize}
\end{frame}

\begin{frame}{Introducción a MapReduce}
\begin{itemize}
  \item \textbf{Map}: transformar cada elemento en pares clave--valor (o en una colección de pares).
  \item \textbf{Shuffle/Sort}: agrupar por clave (fase costosa).
  \item \textbf{Reduce}: combinar todos los valores asociados a cada clave (agregación).
  \item Patrón típico: \textbf{contar}, \textbf{sumar}, \textbf{agrupar}, \textbf{calcular métricas por categoría}.
\end{itemize}
\end{frame}

\begin{frame}{MapReduce en Spark (RDDs)}
\begin{itemize}
  \item Spark implementa este patrón con transformaciones sobre RDDs:
  \begin{itemize}
    \item \texttt{map}/\texttt{flatMap} para generar claves y valores.
    \item \texttt{reduceByKey} o \texttt{aggregateByKey} para reducir por clave.
  \end{itemize}
  \item Diferencia clave: Spark puede \textbf{mantener datos en memoria} y encadenar múltiples etapas de forma eficiente.
  \item \textbf{Buenas prácticas}:
  \begin{itemize}
    \item Preferir \texttt{reduceByKey} frente a \texttt{groupByKey} (menos datos en shuffle).
    \item Controlar el particionado (\texttt{repartition}/\texttt{coalesce}) si hay \emph{skew}.
  \end{itemize}
\end{itemize}
\end{frame}

\begin{frame}[fragile]{Ejemplo: WordCount (MapReduce clásico) con RDD (Scala)}
\begin{verbatim}
val lines = sc.textFile("/data/book.txt")

val counts = lines
  .flatMap(_.split("\\s+"))
  .map(word => (word.toLowerCase, 1))
  .reduceByKey(_ + _)

counts.take(10).foreach(println)
\end{verbatim}
\end{frame}

\begin{frame}{Aplicación práctica: caso típico de logs}
\begin{itemize}
  \item Objetivo: \textbf{nº de eventos por tipo} y \textbf{top N} más frecuentes.
  \item Flujo típico:
  \begin{enumerate}
    \item Leer logs (texto/JSON) \(\rightarrow\) parseo.
    \item Mapear a \((tipoEvento, 1)\).
    \item \texttt{reduceByKey} para contar.
    \item Ordenar y seleccionar top N.
  \end{enumerate}
  \item Métricas a vigilar: particiones, shuffle, tiempo por stage, \emph{skew}.
\end{itemize}
\end{frame}

\begin{frame}[fragile]{Ejemplo práctico: top 5 eventos (Scala)}
\begin{verbatim}
val rdd = sc.textFile("/data/events.log")

val top5 = rdd
  .map(parseEventType)          // String => eventType
  .map(t => (t, 1))
  .reduceByKey(_ + _)
  .map { case (t, c) => (c, t) }
  .sortByKey(ascending = false)
  .take(5)

top5.foreach(println)
\end{verbatim}
\end{frame}

% ------------------------------------------------------------
% Sección
% ------------------------------------------------------------
\section{Iniciación a Spark en cluster}

\begin{frame}{Arquitectura básica de Spark en cluster}
\begin{itemize}
  \item \textbf{Driver}: proceso principal que contiene el \texttt{SparkSession} y construye el DAG.
  \item \textbf{Cluster Manager}: gestiona recursos (CPU, memoria) en los nodos.
  \item \textbf{Executors}: procesos en los nodos worker que ejecutan tareas.
  \item \textbf{Tasks}: unidad mínima de ejecución, una por partición.
  \item Flujo general:
  \begin{enumerate}
    \item El driver solicita recursos al cluster manager.
    \item Se lanzan executors.
    \item El driver envía tareas a los executors.
    \item Los executors procesan particiones y devuelven resultados.
  \end{enumerate}
\end{itemize}
\end{frame}

\begin{frame}{Gestores de recursos soportados}
\begin{itemize}
  \item \textbf{Spark Standalone}
  \begin{itemize}
    \item Cluster manager propio de Spark.
    \item Sencillo de configurar.
    \item Adecuado para entornos dedicados.
  \end{itemize}
  \item \textbf{Hadoop YARN}
  \begin{itemize}
    \item Integrado en ecosistemas Hadoop.
    \item Comparte recursos con otras aplicaciones.
    \item Muy habitual en entornos empresariales.
  \end{itemize}
  \item \textbf{Apache Mesos}
  \begin{itemize}
    \item Cluster manager generalista.
    \item Permite compartir recursos entre múltiples frameworks.
  \end{itemize}
\end{itemize}
\end{frame}

\begin{frame}{Funcionamiento del Driver de Spark}
\begin{itemize}
  \item Punto de entrada de la aplicación.
  \item Construye el \textbf{DAG} de transformaciones.
  \item Divide el job en \textbf{stages} separados por shuffles.
  \item Planifica tareas y las asigna a executors.
  \item Mantiene metadatos, estado de ejecución y resultados parciales.
  \item Si el driver falla, la aplicación termina (salvo configuraciones específicas en cluster).
\end{itemize}
\end{frame}

\begin{frame}{Modos de despliegue del Driver}
\begin{itemize}
  \item \textbf{Client mode}:
  \begin{itemize}
    \item El driver se ejecuta fuera del cluster (por ejemplo, en la máquina local).
    \item Útil para desarrollo y debugging.
  \end{itemize}
  \item \textbf{Cluster mode}:
  \begin{itemize}
    \item El driver se ejecuta dentro del cluster.
    \item Más robusto para producción.
  \end{itemize}
\end{itemize}
\end{frame}

\begin{frame}{Principales propiedades de configuración}
\begin{itemize}
  \item \texttt{spark.master}: URL del cluster manager.
  \item \texttt{spark.app.name}: nombre de la aplicación.
  \item \texttt{spark.executor.memory}: memoria por executor.
  \item \texttt{spark.executor.cores}: núcleos por executor.
  \item \texttt{spark.driver.memory}: memoria del driver.
  \item \texttt{spark.executor.instances}: número de executors.
  \item \texttt{spark.default.parallelism}: paralelismo por defecto.
\end{itemize}
\end{frame}

\begin{frame}[fragile]{Ejemplo de envío de aplicación al cluster}
\begin{verbatim}
spark-submit \
  --class com.ejemplo.Main \
  --master yarn \
  --deploy-mode cluster \
  --executor-memory 4G \
  --executor-cores 2 \
  --num-executors 4 \
  app.jar
\end{verbatim}
\end{frame}

\begin{frame}{Compilación y empaquetamiento}
\begin{itemize}
  \item Aplicaciones Scala/Java:
  \begin{itemize}
    \item Compilación con Maven o SBT.
    \item Generación de JAR con dependencias (fat jar o shaded jar).
  \end{itemize}
  \item Dependencias externas:
  \begin{itemize}
    \item Marcar Spark como \texttt{provided} si el cluster ya lo incluye.
  \end{itemize}
  \item Python:
  \begin{itemize}
    \item Scripts enviados con \texttt{spark-submit}.
    \item Posible uso de entornos virtuales o \texttt{--py-files}.
  \end{itemize}
\end{itemize}
\end{frame}

\begin{frame}{Logging en aplicaciones Spark}
\begin{itemize}
  \item Spark utiliza \textbf{log4j / log4j2} para el sistema de logs.
  \item Permite configurar:
  \begin{itemize}
    \item Nivel de log (INFO, WARN, ERROR, DEBUG).
    \item Appenders (consola, fichero).
  \end{itemize}
  \item Importante para:
  \begin{itemize}
    \item Diagnóstico de errores.
    \item Análisis de rendimiento.
    \item Seguimiento en producción.
  \end{itemize}
  \item Complemento visual: Spark UI (puerto 4040 en local).
\end{itemize}
\end{frame}

\begin{frame}{Resumen conceptual}
\begin{itemize}
  \item El driver orquesta; los executors ejecutan.
  \item El cluster manager asigna recursos.
  \item La configuración adecuada impacta directamente en rendimiento y estabilidad.
  \item El empaquetado y logging correctos son esenciales para producción.
\end{itemize}
\end{frame}

% ------------------------------------------------------------
% Sección
% ------------------------------------------------------------
\section{Conceptos de programación paralela en Spark}

\begin{frame}{Paralelismo en Spark: idea fundamental}
\begin{itemize}
  \item Spark paraleliza el procesamiento a nivel de \textbf{particiones}.
  \item Cada partición se procesa mediante una \textbf{task}.
  \item El número de particiones influye directamente en:
  \begin{itemize}
    \item Nivel de paralelismo.
    \item Uso de CPU.
    \item Overhead de planificación.
  \end{itemize}
  \item Regla práctica: más particiones $\neq$ siempre mejor rendimiento.
\end{itemize}
\end{frame}

\begin{frame}{Particionado de los RDDs cargados desde fichero}
\begin{itemize}
  \item Cuando se usa \texttt{sc.textFile()} o similares:
  \begin{itemize}
    \item Spark crea una partición por \textbf{bloque de fichero} (por defecto).
  \end{itemize}
  \item Puede indicarse número mínimo de particiones:
  \begin{itemize}
    \item \texttt{sc.textFile(path, minPartitions)}
  \end{itemize}
  \item Influye:
  \begin{itemize}
    \item Tamaño del fichero.
    \item Tamaño del bloque en el sistema de almacenamiento.
  \end{itemize}
  \item Buen diseño de particionado mejora el paralelismo inicial.
\end{itemize}
\end{frame}

\begin{frame}[fragile]{Ejemplo: controlando particiones}
\begin{verbatim}
val rdd = sc.textFile("/data/logs.txt", 8)

println(rdd.getNumPartitions)

// Reparticionar después
val repartitioned = rdd.repartition(16)
\end{verbatim}
\end{frame}

\begin{frame}{Particionado de la información leída desde HDFS}
\begin{itemize}
  \item HDFS divide los ficheros en \textbf{bloques} (por ejemplo, 128MB).
  \item Cada bloque suele corresponder a una partición inicial en Spark.
  \item Ventaja clave: \textbf{data locality}.
  \begin{itemize}
    \item Las tareas se intentan ejecutar en el nodo que contiene el bloque.
    \item Reduce tráfico de red.
  \end{itemize}
  \item El paralelismo inicial depende del número de bloques.
\end{itemize}
\end{frame}

\begin{frame}{Repartition vs Coalesce}
\begin{itemize}
  \item \textbf{repartition(n)}:
  \begin{itemize}
    \item Aumenta o disminuye particiones.
    \item Siempre implica shuffle.
  \end{itemize}
  \item \textbf{coalesce(n)}:
  \begin{itemize}
    \item Reduce particiones.
    \item Puede evitar shuffle si se usa correctamente.
  \end{itemize}
  \item Uso típico:
  \begin{itemize}
    \item Reducir particiones antes de escribir a disco.
  \end{itemize}
\end{itemize}
\end{frame}

\begin{frame}{Cálculo de los Stages en Spark}
\begin{itemize}
  \item Spark construye un \textbf{DAG} de transformaciones.
  \item Divide el job en \textbf{stages} separados por operaciones que requieren shuffle.
  \item Regla general:
  \begin{itemize}
    \item Transformaciones \textbf{narrow} $\rightarrow$ mismo stage.
    \item Transformaciones \textbf{wide} (shuffle) $\rightarrow$ nuevo stage.
  \end{itemize}
  \item Cada stage se compone de múltiples \textbf{tasks}.
\end{itemize}
\end{frame}

\begin{frame}{Ejemplo conceptual de división en stages}
\begin{itemize}
  \item Código:
  \begin{enumerate}
    \item \texttt{map}
    \item \texttt{filter}
    \item \texttt{reduceByKey}
    \item \texttt{map}
  \end{enumerate}
  \item Resultado:
  \begin{itemize}
    \item Stage 1: \texttt{map} + \texttt{filter}
    \item Stage 2: \texttt{reduceByKey} (shuffle)
    \item Stage 3: \texttt{map} posterior
  \end{itemize}
\end{itemize}
\end{frame}

\begin{frame}{Relación entre particiones y stages}
\begin{itemize}
  \item Cada stage genera tantas tasks como particiones tenga el RDD de entrada.
  \item Más particiones $\rightarrow$ más tasks $\rightarrow$ mayor paralelismo.
  \item Pero también:
  \begin{itemize}
    \item Más overhead.
    \item Más coste de coordinación.
  \end{itemize}
  \item El equilibrio es clave en programación paralela.
\end{itemize}
\end{frame}

\begin{frame}{Ideas clave para programación paralela eficiente}
\begin{itemize}
  \item Diseñar bien el particionado inicial.
  \item Minimizar shuffles innecesarios.
  \item Evitar \emph{data skew}.
  \item Ajustar número de particiones según recursos del cluster.
  \item Usar Spark UI para analizar stages y tareas.
\end{itemize}
\end{frame}

% ------------------------------------------------------------
% Sección
% ------------------------------------------------------------
\section{Optimización de rendimiento en Spark}

\begin{frame}{El linaje (Lineage) en Spark}
\begin{itemize}
  \item Cada RDD mantiene un \textbf{linaje de transformaciones} que describe cómo fue construido.
  \item Este linaje forma un \textbf{DAG} (Directed Acyclic Graph).
  \item Permite:
  \begin{itemize}
    \item \textbf{Tolerancia a fallos}: recomputar particiones perdidas.
    \item Optimización de ejecución.
  \end{itemize}
  \item Implicación práctica:
  \begin{itemize}
    \item Si no se persiste un RDD, puede recomputarse varias veces.
  \end{itemize}
\end{itemize}
\end{frame}

\begin{frame}{Coste del linaje largo}
\begin{itemize}
  \item Pipelines muy largos $\rightarrow$ DAG complejo.
  \item Si hay múltiples acciones sobre el mismo RDD:
  \begin{itemize}
    \item Se recalcula desde el origen si no está persistido.
  \end{itemize}
  \item Solución:
  \begin{itemize}
    \item \textbf{cache()} o \textbf{persist()}.
    \item \textbf{checkpoint()} para cortar el linaje en procesos muy largos.
  \end{itemize}
\end{itemize}
\end{frame}

\begin{frame}{Estrategias de almacenamiento de RDDs}
\begin{itemize}
  \item Spark permite distintos niveles de persistencia:
  \begin{itemize}
    \item \texttt{MEMORY\_ONLY}
    \item \texttt{MEMORY\_AND\_DISK}
    \item \texttt{DISK\_ONLY}
    \item Versiones serializadas
  \end{itemize}
  \item Elección depende de:
  \begin{itemize}
    \item Tamaño del dataset.
    \item Memoria disponible.
    \item Coste de recomputación.
  \end{itemize}
\end{itemize}
\end{frame}

\begin{frame}[fragile]{Ejemplo de persistencia}
\begin{verbatim}
val rdd = sc.textFile("/data/logs.txt")
  .map(parse)
  .filter(valid)

rdd.cache()

val count1 = rdd.count()
val count2 = rdd.filter(_.level == "ERROR").count()
\end{verbatim}
\end{frame}

\begin{frame}{Buenas prácticas de persistencia}
\begin{itemize}
  \item Persistir solo si:
  \begin{itemize}
    \item Se reutiliza varias veces.
    \item La recomputación es costosa.
  \end{itemize}
  \item Liberar memoria cuando ya no se necesita:
  \begin{itemize}
    \item \texttt{rdd.unpersist()}
  \end{itemize}
  \item Evitar cachear datasets enormes sin memoria suficiente.
\end{itemize}
\end{frame}

\begin{frame}{Variables de Broadcast}
\begin{itemize}
  \item Permiten enviar una copia \textbf{eficiente y distribuida} de un objeto a todos los executors.
  \item Evitan:
  \begin{itemize}
    \item Enviar el mismo objeto en cada task.
    \item Sobrecarga de red innecesaria.
  \end{itemize}
  \item Uso típico:
  \begin{itemize}
    \item Tablas pequeñas de referencia.
    \item Diccionarios de lookup.
  \end{itemize}
\end{itemize}
\end{frame}

\begin{frame}[fragile]{Ejemplo de Broadcast}
\begin{verbatim}
val lookup = Map("ES" -> "Spain", "FR" -> "France")

val broadcastVar = sc.broadcast(lookup)

val enriched = rdd.map { row =>
  val country = broadcastVar.value.get(row.code)
  (row, country)
}
\end{verbatim}
\end{frame}

\begin{frame}{Acumuladores (Accumulators)}
\begin{itemize}
  \item Variables que permiten agregar información desde los executors al driver.
  \item Uso típico:
  \begin{itemize}
    \item Contadores.
    \item Métricas de calidad.
  \end{itemize}
  \item Características:
  \begin{itemize}
    \item Solo el driver puede leer el valor final.
    \item Los executors solo pueden incrementar.
  \end{itemize}
\end{itemize}
\end{frame}

\begin{frame}[fragile]{Ejemplo de Acumulador}
\begin{verbatim}
val errorCounter = sc.longAccumulator("ErrorCounter")

rdd.foreach { row =>
  if (!isValid(row)) {
    errorCounter.add(1)
  }
}

println(errorCounter.value)
\end{verbatim}
\end{frame}

\begin{frame}{Mecanismos simples para mejorar rendimiento}
\begin{itemize}
  \item Minimizar shuffles.
  \item Ajustar número de particiones.
  \item Persistir correctamente.
  \item Usar broadcast en joins pequeños.
  \item Utilizar acumuladores para métricas sin romper paralelismo.
  \item Monitorizar con Spark UI.
\end{itemize}
\end{frame}

% ------------------------------------------------------------
% Sección
% ------------------------------------------------------------
\section{Iniciación al procesamiento de datos en Spark mediante el uso del módulo SparkSQL}

\begin{frame}{¿Qué es SparkSQL?}
\begin{itemize}
  \item Módulo de Spark para procesamiento estructurado.
  \item Introduce el concepto de \textbf{DataFrame}.
  \item Basado en optimización automática mediante:
  \begin{itemize}
    \item \textbf{Catalyst Optimizer}
    \item \textbf{Tungsten Execution Engine}
  \end{itemize}
  \item Permite trabajar con:
  \begin{itemize}
    \item API en Scala/Python
    \item SQL directamente
  \end{itemize}
\end{itemize}
\end{frame}

\begin{frame}{Construcción de DataFrames en Spark}
\begin{itemize}
  \item Desde ficheros:
  \begin{itemize}
    \item \texttt{read.csv()}
    \item \texttt{read.json()}
    \item \texttt{read.parquet()}
  \end{itemize}
  \item Desde RDDs.
  \item Desde colecciones locales.
  \item Desde tablas externas (Hive, JDBC).
\end{itemize}
\end{frame}

\begin{frame}[fragile]{Ejemplo: Crear DataFrame desde fichero}
\begin{verbatim}
val df = spark.read
  .option("header", "true")
  .option("inferSchema", "true")
  .csv("/data/users.csv")

df.printSchema()
\end{verbatim}
\end{frame}

\begin{frame}{Transformaciones comunes sobre DataFrames}
\begin{itemize}
  \item \textbf{select}: seleccionar columnas.
  \item \textbf{selectExpr}: expresiones SQL.
  \item \textbf{withColumn}: añadir/modificar columnas.
  \item \textbf{where / filter}: filtrar filas.
  \item \textbf{join}: combinar DataFrames.
  \item \textbf{unionAll}: unir datasets verticalmente.
  \item \textbf{groupBy}: agregaciones.
  \item \textbf{orderBy}: ordenar resultados.
\end{itemize}
\end{frame}

\begin{frame}[fragile]{Ejemplo de transformaciones}
\begin{verbatim}
import org.apache.spark.sql.functions._

val transformed = df
  .select("name", "age", "country")
  .withColumn("age_plus_1", col("age") + 1)
  .filter(col("age") > 18)
  .groupBy("country")
  .count()
  .orderBy(desc("count"))
\end{verbatim}
\end{frame}

\begin{frame}{Tipos de joins en SparkSQL}
\begin{itemize}
  \item \textbf{inner join}
  \item \textbf{left / right outer join}
  \item \textbf{full outer join}
  \item \textbf{cross join}
  \item \textbf{broadcast join} (optimización automática o explícita)
\end{itemize}
\end{frame}

\begin{frame}{Acciones comunes sobre DataFrames}
\begin{itemize}
  \item \textbf{show}: visualizar resultados.
  \item \textbf{collect}: traer datos al driver.
  \item \textbf{count}: contar filas.
  \item \textbf{first / head}: primera fila.
  \item \textbf{take}: primeras N filas.
  \item \textbf{write}: guardar en almacenamiento.
\end{itemize}
\end{frame}

\begin{frame}[fragile]{Ejemplo de acciones}
\begin{verbatim}
transformed.show()

val total = transformed.count()

val firstRow = transformed.first()

transformed.write
  .mode("overwrite")
  .parquet("/data/output")
\end{verbatim}
\end{frame}

\begin{frame}{Buenas prácticas en SparkSQL}
\begin{itemize}
  \item Preferir DataFrames frente a RDDs en procesamiento estructurado.
  \item Evitar \texttt{collect()} en datasets grandes.
  \item Aprovechar funciones built-in antes que UDFs.
  \item Revisar el plan con:
  \begin{itemize}
    \item \texttt{df.explain(true)}
  \end{itemize}
  \item Monitorizar ejecución en Spark UI.
\end{itemize}
\end{frame}

\end{document}